\documentclass[11pt,a4paper]{article}
\usepackage[T1]{fontenc}
\usepackage[utf8]{inputenc}
\usepackage{lmodern}
\usepackage{amsmath,amssymb,amsthm,mathtools,mathrsfs}
\usepackage{graphicx,float,xcolor,multicol}
\usepackage[shortlabels]{enumitem}
\usepackage[margin=27mm]{geometry}
\usepackage{microtype,needspace,placeins}
\usepackage{tikz,tikz-cd}
\usetikzlibrary{arrows.meta,calc,positioning,patterns,decorations.pathreplacing}
\usepackage[colorlinks=true,linkcolor=teal,urlcolor=teal,citecolor=teal]{hyperref}
\setlength{\parindent}{0pt}
\setlength{\parskip}{.6em}
\setcounter{secnumdepth}{0}
\allowdisplaybreaks
\raggedbottom
\providecommand{\cross}{\times}
\DeclareUnicodeCharacter{2020}{\ensuremath{\dagger}}
\DeclareMathOperator{\supp}{supp}
\DeclareMathOperator*{\esssup}{ess\,sup}
\providecommand{\chapter}{\section}
\newenvironment{tool}[2]{\subsection*{Tool #1 --- #2}}{}
\newcommand{\ToolRef}[1]{Tool #1}
\newcommand{\FigureTag}[2]{\textit{Figure #1}}
\newcommand{\Result}[1]{\subsection*{#1}}
\newcommand{\Application}[1]{\subsubsection*{Application\if\relax\detokenize{#1}\relax\else\space--- #1\fi}}
\newcommand{\ProofHeading}{\par\textit{Proof.}\space}
\newcommand{\ProofOf}[1]{\par\textit{Proof of #1.}\space}
\newcommand{\RemarkHeading}{\par\textbf{Remark.}\space}
\newcommand{\NoteHeading}{\par\textbf{Note.}\space}
\newcommand{\ExerciseHeading}{\par\textbf{Exercise.}\space}

\graphicspath{{figures/k-theory/}}
\begin{document}
\section*{Bott Periodicity and the Genesis of K-Theory}
% BLOG-CONTENT-BEGIN
As $K$ is represented by the space $BU\times \mathbb{Z}$, any theorem about the homotopy type of $BU\times \mathbb{Z}$ can be significant. Bott periodicity gives the key homotopy equivalence. 
\begin{theorem}[Bott periodicity]
$BU\times \mathbb{Z} \simeq \Omega^2(BU\times \mathbb{Z})$.
\end{theorem}

For a proof, see \cite{milnor2016morse}. An immediate consequence of Bott's theorem is the following theorem which we will also call Bott's Periodicity Theorem.
\begin{theorem}[Bott periodicity for reduced K-theory]
    $\widetilde{K}(\Sigma^2X) \cong \widetilde{K}(X)$
\end{theorem}

\begin{proof}
    Observe that,
    $$\widetilde{K}(X) \cong [(X, x_0), (BU\times \mathbb{Z}, (y_0, 0))]$$
    as it is the kernel of the map $[X, BU\times \mathbb{Z}] \rightarrow [x_0, BU\times \mathbb{Z}]$, which is precisely the set of homotopy classes of maps sending $x_0$ into $BU\times \{0\}$. Now, by Bott's Theorem,
    $$\widetilde{K}(X) \cong [X, \Omega^2(BU\times \mathbb{Z})]_{\text{based}} \cong [\Sigma^2X, BU\times \mathbb{Z}]_{\text{based}}$$
    giving us, 
    $$\widetilde{K}(X) \cong \widetilde{K}(\Sigma^2X)$$
\end{proof}


An important consequence of the theorem above is that one can build a cohomology theory on $\tau_c$. Generally speaking, a \textbf{cohomology theory} on a category $\mathcal{C}$ of topological pairs is the following data:
$$F^k: \mathcal{C}^{\text{op}} \rightarrow \text{Ab (category of abelian groups)}$$
$\{F^k\}_{k}$ a sequence of presheaves on $\mathcal{C}$ with natural transformations,
$$\delta: F^k(A) := F^k(A, \emptyset) \rightarrow F^{k+1}(X, A)$$
satisfying the following axioms \cite{eilenberg1945axiomatic}:

\underline{Axiom 1} (Homotopy Axiom): If $f$, $g$ $:(X, A) \rightarrow (Y,B)$ are homotopic as maps of pairs, then $f^*, g^*$ $:F^k(Y, B) \rightarrow F^k(X, A)$ are equal, $\forall k\in \mathbb{Z}$, where $f^*:= F^k(f)$.

\underline{Axiom 2} (Sequence Axiom): For a pair $(X, A)$, the above maps induce a long exact sequence:
$$\cdots \rightarrow F^k(X, A)\rightarrow F^k(X) \rightarrow F^k(A) \rightarrow F^{k+1}(X, A)\rightarrow \cdots$$


\underline{Axiom 3} (Excision Axiom): If $U\subseteq \overline{U} \subset \text{Int}(A)$, then $F^k(X\backslash U, A\backslash U) \cong F^k(X, A)$, $\forall k$, $k\in \mathbb{Z}$.

The construction of K-theory on pairs of compact Hausdorff spaces uses the following two propositions \cite{hatcher2003vector}.

\begin{proposition}
Consider $A\xrightarrow{i} X \xrightarrow{q} X/A$, where $i$, $q$ are the canonical inclusion and quotient maps. Then,
$$\widetilde{K}(X/A) \xrightarrow{q^*} \widetilde{K}(X) \xrightarrow{i^*} \widetilde{K}(A)$$
is an exact sequence of rings.
\end{proposition}

\begin{proof}
    Note that $i^*\circ q^* = (q\circ i)^*$, where $q\circ i$ is the constant map from $A$ to the basepoint of $X/A$. So, $\text{Im}(q^*) \subseteq \text{Ker}(i^*)$. Now we show $\text{Ker}(i^*)  \subseteq \text{Im}(q^*)$. If $[E] \in \text{Ker}(i^*)$, $\big[E\big|_{A}\big] = 0$, in other words,
    $$E\big|_{A}\oplus \epsilon^n(A) \cong \epsilon^{n+m}(A)$$
    Therefore, we may WLOG assume that $E\big|_{A}$ is trivial, with trivialisation,
    $$\psi: E\big|_{A} \rightarrow A\times \mathbb{C}^n$$
    Let $E/\psi$ denote the quotient space, where $\psi^{-1}(x, v) \sim \psi^{-1}(y, v)$, for all $x, y\in A$. We claim that $E/\psi$ is a bundle over $X/A$, and that $q^*(E/\psi) \cong E$. To show the former, note that for a locally trivializing cover $\{U_{\alpha}\}_{\alpha}$ of $E$, the modified collection of open\footnote{The spaces $A$, $X$ being compact Hausdorff ensure that the modified collection is indeed a collection of open sets. Perhaps, this dependence of the compactness hypothesis in this proof may be lifted.} sets $\{U_{\alpha}\backslash A\}_{\alpha}$ defines a bundle over $X\backslash A$. This bundle is a subbundle of $E$, and for that to extend to a bundle on $X/A$, one would need a local trivialization at the point $q\circ i(A)$. As $E\big|_{A}$ is trivial, one has a frame $\sigma_i$ over $A$ which can (by a partition of unity argument) be extended to a frame over an open set containing $A$, since the set of points $p$ satisfying
    $$\det(\sigma_1(p), \sigma_2(p),\ldots, \sigma_n(p)) \neq 0$$
    is open, thereby facilitating the required extension. Now note that by definition, the natural map $E\rightarrow E/\psi$ lifts the quotient map, (i.e.) one has the following commutative diagram:
    \[\begin{tikzcd}[column sep=huge,row sep=large]
	E & {E/\psi} \\
	X & {X/A}
	\arrow["q"', from=2-1, to=2-2]
	\arrow[from=1-1, to=2-1]
	\arrow[from=1-2, to=2-2]
	\arrow[from=1-1, to=1-2]
\end{tikzcd}\]
which is also an isomorphism on the fibres. It is easy to see that $E$ satisfies the universal property of pullbacks.
\end{proof}

\begin{remark}
Interpreting the proposition in terms of the classifying space perspective yields the following: Given a classifying map $f$ of some bundle over $X$, such that $f\big|_A$ is nullhomotopic, then there exists an open neighbourhood $U$ of
$A$ such that $f\big|_{U}$ is nullhomotopic. Equivalently, maps from compact subsets into $BU$ that are null-homotopic extend as null-homotopic maps over a neighbourhood. If you think about it, the above is not unexpected at all, for the spaces $BU(n)$ are locally contractible. 
\end{remark}

\begin{proposition}
If $(X,A)$ is a compact Hausdorff pair, the inclusion $A\hookrightarrow X$ is a cofibration, and $A$ is contractible, then $q^*: \widetilde{K}(X/A) \rightarrow \widetilde{K}(X)$ is an isomorphism.
\end{proposition}

\begin{proof}
    We begin by identifying $\widetilde{K}(X)$ with $\pi_b(X,BU)$. Following the remark above, note that if $A$ is contractible then for any classifying map $f: X\rightarrow BU$, $f\big|_A$ is nullhomotopic, and that there exists an open neighbourhood $U$ of $A$ such that $f\big|_U$ is also nullhomotopic. Let,
    $$F:U\times [0,1] \rightarrow BU$$
    $$F_0\equiv f\big|_U, F_1 \equiv b_0$$ be the nullhomotopy. As compact Hausdorff spaces are normal, by Urysohn's lemma we have, $$\alpha :X \rightarrow [0,1]$$
    $$\alpha\big|_A = 1, \supp{\alpha} \subset U$$
    Define $Q(f): X \rightarrow BU$ by,

        \begin{equation*}
            Q(f) := 
            \begin{cases}
                F(x,\alpha(x)),& \text{ if } x\in U\\
            f(x), & \text{ else}
            \end{cases}
        \end{equation*}
        which is clearly continuous. As $Q(f)$ is constant on $A$, it factors through $X/A \rightarrow BU$. It is easy to see that  the induced maps give the inverse of $q^*$, so we call them $(q^*)^{-1}(f)$, for $f\in \pi_b(X, BU)$. It remains to check that the map $Q$ is well-defined, which essentially boils down to proving that $Q(f)$ depends only on the homotopy type of $f$. Let $f$, $g$ $:X \rightarrow BU$ be homotopic, with,
        $$H: X\times [0,1] \rightarrow BU$$
        $$H_0 \equiv f, H_1 \equiv g$$
        the homotopy between them. Now as $A$ is contractible, so is $A\times [0,1]$, implying that, $H\big|_{A\times [0,1]}$ is nullhomotopic. Let the homotopy be given by the map $N$. Now, again apply Urysohn's lemma for the space $X\times [0,1]$ to get,
        $$\beta: X\times [0,1] \rightarrow [0,1]$$
        $$\beta\big|_{A\times [0,1]} \equiv 1, \supp{\beta} \subset V\times [0,1]$$
        Repeat the construction above by defining, 
        $$G(H_t):= \begin{cases}
            N((x, t), \beta(x,t)), & \text{ if } x\in V\\
            H_t(x), & \text{ else}
        \end{cases}$$
        To finish the proof, let's do away with one small yet essential detail: One can define a straight-line homotopy between $\alpha$ and $\beta_t$, hence the function $Q$ is independent of the open neighbourhoods on which $f$ and $g$ are null-homotopic. 
\end{proof}

Let $n \in \mathbb{N}\cup \{0\}$, define,
$$\widetilde{K}^{-n}(X, A):= \widetilde{K}(\Sigma^n(X/A)).$$
Using Bott periodicity, extend this definition to all integer degrees by $\widetilde K^{m+2}(X,A)\cong\widetilde K^m(X,A)$.
For the sequence of functors defined above, homotopy axiom follows from the representation of the functor $K$, and the excision axiom is just the trivial statement $(X\backslash U)/(A\backslash U) = X/A$. To investigate the sequence axiom, we need to define the requisite natural transformations; for this purpose, Proposition $4.2$ will come in handy. Given a pair $(X, A)$ we have the following sequence of spaces called the \textbf{Puppe sequence}:
\[\begin{tikzcd}
	&& {(((X\cup CA)\cup CX)\cup C(X\cup CA))/C(X\cup A))\sim \color{red}{\Sigma} X} \\
	A & \cdots & {((X\cup CA)\cup CX)\cup C(X\cup CA)} \\
	X & {X\cup CA} & {(X\cup CA)\cup CX} \\
	& {(X\cup CA)/CA \sim \color{red}{X/A}} & {((X\cup CA)\cup CX)/CX \sim \color{red}{\Sigma A}}
	\arrow[hook, from=3-1, to=3-2]
	\arrow[hook, from=3-2, to=3-3]
	\arrow[two heads, from=3-2, to=4-2]
	\arrow[hook', from=2-3, to=2-2]
	\arrow[hook, from=2-1, to=3-1]
	\arrow[two heads, from=2-3, to=1-3]
	\arrow[two heads, from=3-3, to=4-3]
\end{tikzcd}\]


Note that as cones are contractible, we have the isomorphism: 
$$\widetilde{K}(X\cup CA) \cong \widetilde{K}((X\cup CA)/CA) \cong \widetilde{K}(X/A)$$
which gives rise to a cofibration-like sequence:
$$\widetilde{K}(A) \xleftarrow{} \widetilde{K}(X) \xleftarrow{} \widetilde{K}(X/A) \xleftarrow{} \widetilde{K}(\Sigma A) \xleftarrow{} \widetilde{K}(\Sigma X) \xleftarrow{} \widetilde{K}(\Sigma X/\Sigma A) \xleftarrow{} \cdots$$
Therefore by Bott's Theorem, the above induces the following:
\[\begin{tikzcd}
	&&&&& {\widetilde{K}^{1}(X/A)} & {\widetilde{K}^{1}(X)} \\
	\cdots & {\widetilde{K}^0(X/A)} & {\widetilde{K}^0(X)} & {\widetilde{K}^{0}(A)} & {\widetilde{K}^{-2}(A)} & {\widetilde{K}^{-1}(X/A)} & {\widetilde{K}^{-1}(X)} \\
	&& \cdots & {\widetilde{K}^{0}(A)} & {\widetilde{K}^{0}(X)} & {\widetilde{K}^{0}(X/A)} & {\widetilde{K}^{-1}(A)} \\
	&&& {\widetilde{K}^{2}(A)} & {\widetilde{K}^{2}(X)} & {\widetilde{K}^{2}(X/A)}
	\arrow[from=2-1, to=2-2]
	\arrow[from=2-2, to=2-3]
	\arrow[from=2-3, to=2-4]
	\arrow["\cong", from=2-4, to=2-5]
	\arrow[from=2-5, to=2-6]
	\arrow[from=2-6, to=2-7]
	\arrow[from=2-7, to=3-7]
	\arrow[from=3-7, to=3-6]
	\arrow[from=3-6, to=3-5]
	\arrow[from=3-5, to=3-4]
	\arrow["\cong"', from=3-4, to=4-4]
	\arrow["\cong"', from=3-5, to=4-5]
	\arrow["\cong", from=2-7, to=1-7]
	\arrow["\cong", from=2-6, to=1-6]
	\arrow["\cong"', from=3-6, to=4-6]
	\arrow[from=3-4, to=3-3]
\end{tikzcd}\]
which, written in a more recognisable form, is 
$$\cdots \rightarrow \widetilde{K}^n(X, A)\rightarrow \widetilde{K}^n(X) \rightarrow \widetilde{K}^n(A) \rightarrow \widetilde{K}^{n+1}(X, A)\rightarrow \cdots$$

The periodicity theorem essentially allows us to  wrap the sequence on itself. Exactness at various places will be given by Proposition $4.1$ applied to relevant sections of the Puppe sequence, along with the observation that $(X\cup CY)/CY \sim X/Y$. So, $\{K^n\}_{n}$ satisfies the sequence axiom for pairs. Thus $\{\widetilde K^n\}_n$ satisfies the axioms of a generalised cohomology theory on compact Hausdorff pairs.
\begin{theorem}
    The sequence of functors $\widetilde{K}^n$ defines a generalized cohomology theory on the category of compact Hausdorff pairs.
\end{theorem}

We finish by computing the $K$-groups of spheres and tori to illustrate the theory. Let $X$ and $Y$ be compact Hausdorff; one has the following exact sequence of pointed spaces:
\begin{equation}\tag{2}
    X\vee Y \hookrightarrow X\times Y \rightarrow X\wedge Y:= (X\times Y)/(X\vee Y)
\end{equation}
\begin{equation}\tag{3}
    X\hookrightarrow X\vee Y \rightarrow Y \cong (X\vee Y)/X
\end{equation}
Note that $(2)$ and $(3)$ yield the following split exact sequence:
$$0\rightarrow \widetilde{K}(X\wedge Y) \rightarrow \widetilde{K}(X\times Y) \rightarrow \widetilde{K}(X)\oplus \widetilde{K}(Y) \xrightarrow{} 0$$
which gives the isomorphism,
\begin{equation}\tag{4}
    \widetilde{K}(X\times Y) \cong \widetilde{K}(X)\oplus \widetilde{K}(Y)\oplus \widetilde{K}(X\wedge Y)
\end{equation}


\textbf{Computations}:
\begin{enumerate}[(i)]
    \item (K-Groups of Spheres): Note that, $\widetilde{K}(S^k) \cong \pi_k(BU) \cong \pi_{k+1}(U)$; here we exploit the fact implicit in Bott's Theorem that $U\sim \Omega BU$. Now, we need to compute two distinct homotopy groups of the stable unitary group $U$. Note,
    $$\pi_1(U) = \pi_1(U(1)) = \pi_1(S^1) \cong \mathbb{Z}$$
    The fibration $SU(2)\rightarrow U(2)\rightarrow S^1$ and the identification $SU(2)\cong S^3$ give
    $$\pi_2(U) = \pi_2(U(2)) \cong \pi_2(SU(2)) \cong \pi_2(S^3) = 0.$$
    By Bott periodicity, we have $\pi_{2n}(U) \cong 0$, and $\pi_{2n+1}(U) \cong \mathbb{Z}$, for $n\in \mathbb{N}$. This gives us that:
    $$\widetilde{K}(S^n) = \begin{cases}
        \mathbb{Z}, &\text{ if }2|n\\
        0, & \text{ else}
    \end{cases}$$
    \item (K-Groups of $n$-Torus): Apply $(4)$ to the spaces $T^{n-1}$ and $S^1$, to get:
    $$\widetilde{K}^i(T^n) \cong \widetilde{K}^i(T^{n-1})\oplus \widetilde{K}^i(S^1)\oplus \widetilde{K}^{i+1}(T^{n-1})$$
    By induction one gets the following:
    $$\widetilde{K}^i(T^n)= \begin{cases}
        \mathbb{Z}^{\oplus(2^{n-1} -1)}, & \text{ if } 2|i\\
        \mathbb{Z}^{\oplus(2^{n-1})}, & \text{ else}
    \end{cases}$$
\end{enumerate}

% BLOG-CONTENT-END
\bibliographystyle{amsalpha}
\bibliography{tools/profiles/k-theory/references}
\end{document}
